% =====================================================================
% dual_formulation.tex --- D4 --- LP-duality derivation of lambda <-> eps.
% Inserted into sections/methodology.tex AFTER the operational-state
% subsection and BEFORE the marginal-emission-factor subsection.
% =====================================================================

\subsection{Dual-LP Formulation: From $\lambda$ to $\varepsilon$}
\label{sec:dual-formulation}

The closest-CV picker introduced in
Section~\ref{sec:scvic-ladder} can be reinterpreted as the empirical
realisation of a chain of linear-programming dualities.  Stating the
chain explicitly serves three purposes: (i)~it pins down the
sufficient conditions under which the reconstructed
$\hat\varepsilon_t$ is provably equal to the true marginal emission
factor of the underlying physical dispatch, (ii)~it isolates the
assumptions that fail when the merit stack is non-strict, and
(iii)~it provides the analytical scaffold for the BESS case study of
Section~\ref{sec:bess-case}, in which a carbon adder
$p_{\mathrm{CO}_2}$ is folded into the arbitrage signal.  The
schematic of the three primal/dual pairs developed below is given in
Fig.~\ref{fig:lp-duality}.

\begin{figure}[t]
  \centering
  % =====================================================================
% lp_duality.tex --- primal/dual LP for cost and emission, with
% lambda <-> epsilon as parallel duals.
% =====================================================================
\begin{tikzpicture}[
  font=\scriptsize,
  box/.style={draw, rounded corners=3pt, inner sep=4pt,
              text width=44mm, align=left,
              font=\scriptsize},
  cbox/.style={box, fill=blue!8},
  ebox/.style={box, fill=green!10},
  jbox/.style={box, fill=orange!10},
  arr/.style={draw, -{Stealth[length=4pt]}, thick},
  darr/.style={draw, -{Stealth[length=4pt]}, thick, dashed},
]

  % cost LP --- primal
  \node[cbox] (cp) at (0,0)
    {\textbf{Cost primal} (LP-c)\\[1pt]
     $\displaystyle\min_p \sum_g c_g\, p_g$\\
     s.t.\ $\sum_g p_g = D_t\;\,[\lambda_t]$\\
     \phantom{s.t.\ }$0\!\le\!p_g\!\le\!\bar p_g\;\,[\mu_g]$};

  % cost LP --- dual
  \node[cbox] (cd) [right=22mm of cp]
    {\textbf{Cost dual} (LP-c$^\star$)\\[1pt]
     $\displaystyle\max_{\lambda,\mu}\; D_t\lambda - \!\sum_g\bar p_g\mu_g$\\
     s.t.\ $\lambda - \mu_g \le c_g\;\forall g$\\
     KKT: $\lambda_t = c_{m(t)}$};

  % emission LP --- primal
  \node[ebox] (ep) [below=10mm of cp]
    {\textbf{Emission primal} (LP-e)\\[1pt]
     $\displaystyle\min_p \sum_g e_g\, p_g$\\
     s.t.\ $\sum_g p_g = D_t\;\,[\varepsilon_t]$\\
     \phantom{s.t.\ }$0\!\le\!p_g\!\le\!\bar p_g\;\,[\nu_g]$};

  % emission LP --- dual
  \node[ebox] (ed) [right=22mm of ep]
    {\textbf{Emission dual} (LP-e$^\star$)\\[1pt]
     $\displaystyle\max_{\varepsilon,\nu}\;
        D_t\varepsilon - \!\sum_g\bar p_g\nu_g$\\
     s.t.\ $\varepsilon - \nu_g \le e_g\;\forall g$\\
     KKT: $\varepsilon_t = e_{m'(t)}$};

  % joint LP
  \node[jbox] (jp) [below=10mm of ep]
    {\textbf{Joint primal} (LP-j, $p_{\mathrm{CO}_2}\!>\!0$)\\[1pt]
     $\displaystyle\min_p \sum_g (c_g\!+\!p_{\mathrm{CO}_2}e_g)p_g$\\
     s.t.\ $\sum_g p_g = D_t\;\,[\tilde\lambda_t]$};

  \node[jbox] (jd) [right=22mm of jp]
    {\textbf{Joint dual} (LP-j$^\star$)\\[1pt]
     KKT: $\tilde\lambda_t =
            c_{m''(t)} + p_{\mathrm{CO}_2}\!\cdot\! e_{m''(t)}$\\
     $\Rightarrow$ $m''(t)$ shifts as
     $p_{\mathrm{CO}_2}\!\uparrow$};

  % strong duality arrows
  \draw[arr] (cp) -- node[above,font=\tiny]
                    {strong duality} (cd);
  \draw[arr] (ep) -- node[above,font=\tiny]
                    {strong duality} (ed);
  \draw[arr] (jp) -- node[above,font=\tiny]
                    {strong duality} (jd);

  % parallel between cost and emission
  \draw[darr] (cd.south) -- node[right,font=\tiny]
                            {analogue (swap $c_g\!\to\!e_g$)} (ed.north);
  \draw[darr] (ed.south) -- node[right,font=\tiny]
                            {add $p_{\mathrm{CO}_2}$ adder} (jd.north);

\end{tikzpicture}

  \caption{Cost, emission, and joint primal/dual LPs.
    Strong duality and the swap $c_g\!\leftrightarrow\!e_g$ produce the
    parallelism that justifies reconstructing $\varepsilon_t$ from a
    published $\lambda_t$ without re-solving an LP.}
  \label{fig:lp-duality}
\end{figure}

\paragraph{(a) Cost-minimising primal and its dual.}
Consider an instantaneous economic dispatch at block $t$, ignoring
transmission for clarity.  The primal LP is
\begin{align}
  \min_{p \in \R^{|\setG|}_{\geq 0}}
    \quad & \sum_{g \in \setG} c_g\, p_g
        \label{eq:cost-primal} \\
  \text{s.t.}\quad
    & \sum_{g \in \setG} p_g \;=\; D_t
        \quad &&[\lambda_t]
        \label{eq:cost-balance}\\
    & p_g \;\le\; \bar p_g
        \quad \forall g\in\setG
        \quad &&[\mu_g\!\geq\!0]
        \label{eq:cost-bound}\\
    & p_g \;\geq\; 0 \;\;\forall g.\nonumber
\end{align}
The associated dual LP is
\begin{align}
  \max_{\lambda,\,\mu\geq 0}
  \quad &  D_t\,\lambda \;-\; \sum_{g} \bar p_g\, \mu_g
  \quad
  \text{s.t.}\quad \lambda - \mu_g \le c_g \;\forall g.
  \label{eq:cost-dual}
\end{align}
Strong duality holds because~\eqref{eq:cost-primal}--\eqref{eq:cost-bound}
is a feasible bounded LP.  At the optimum the
Karush--Kuhn--Tucker conditions yield, for any generator $g$ that is
strictly interior to its bounds ($0<p_g^\star<\bar p_g$),
\begin{equation}
  \mu_g^\star = 0
  \quad\Longrightarrow\quad
  \lambda_t^\star = c_g.
  \label{eq:lambda-equals-cm}
\end{equation}
Denoting by $m(t)$ the index of any such interior unit (the
\emph{cost-marginal} unit), \eqref{eq:lambda-equals-cm} reads
$\lambda_t^\star = c_{m(t)}$.  This is the textbook merit-order
identity~\cite{SchweppeBook1988,Hogan1992}.

\paragraph{(b) Emission-minimising primal and its dual.}
Replacing the cost coefficient $c_g$ by an emission coefficient
$e_g\!=\!\mathrm{EF}_{\mathrm{fuel}(g)}\!\cdot\!\mathrm{HR}_g$
(tCO$_2$/MWh) leaves the constraint structure of
\eqref{eq:cost-balance}--\eqref{eq:cost-bound} intact and produces the
\emph{emission primal}
\begin{align}
  \min_{p}
    \quad & \sum_{g} e_g\, p_g
    \quad
    \text{s.t.}\quad
    \sum_{g} p_g = D_t \;\,[\varepsilon_t],\;\;
    p_g \le \bar p_g \;\,[\nu_g],\;\;
    p_g \ge 0.
        \label{eq:emis-primal}
\end{align}
By the same derivation as in (a), the dual variable on the demand
balance reads
\begin{equation}
  \varepsilon_t^\star = e_{m'(t)},
  \label{eq:eps-equals-em}
\end{equation}
where $m'(t)$ is the \emph{emission-marginal} unit.  In general
$m(t)\neq m'(t)$: when, for example, an inexpensive coal unit and a
more expensive gas unit both lie at the cost margin in different
hours, the cost-marginal stack rotates between coal and gas while
the emission-marginal stack rotates only between fuels of distinct
$e$-coefficient.

\paragraph{(c) Joint primal with a carbon adder.}
Adding a carbon price $p_{\mathrm{CO}_2}$ (USD/tCO$_2$) to every
generator cost converts the cost coefficient into the
\emph{social cost}
$\tilde c_g \;=\; c_g + p_{\mathrm{CO}_2}\,e_g$.
The joint primal is
\begin{align}
  \min_{p}
  \quad & \sum_{g} \big(c_g + p_{\mathrm{CO}_2}\, e_g\big)\, p_g
    \quad
  \text{s.t.}\quad
  \sum_g p_g = D_t \;\,[\tilde\lambda_t],\;\;
  p_g \le \bar p_g.
  \label{eq:joint-primal}
\end{align}
Its dual produces, for any interior generator $g\!=\!m''(t)$,
\begin{equation}
  \tilde\lambda_t^\star
    \;=\; c_{m''(t)} \;+\; p_{\mathrm{CO}_2}\, e_{m''(t)}.
  \label{eq:joint-lambda}
\end{equation}
The marginal unit $m''(t)$ is in general distinct from both $m(t)$
and $m'(t)$, and the substitution between two competing units
$g\!=\!1$ and $g\!=\!2$ (with $c_1<c_2$ but $e_1>e_2$) admits an
analytical \emph{switching carbon price}
\begin{equation}
  p^\star \;=\; \frac{c_2 - c_1}{e_1 - e_2},
  \label{eq:switching-price}
\end{equation}
above which unit~2 displaces unit~1 in the merit order.  For the
representative Chilean coal-versus-gas pair
($c_{\rm coal}\!\approx\!\SI{45}{\dollar/\mega\watt\hour}$,
$e_{\rm coal}\!\approx\!\SI{0.94}{tCO_2/\mega\watt\hour}$;
$c_{\rm gas}\!\approx\!\SI{75}{\dollar/\mega\watt\hour}$,
$e_{\rm gas}\!\approx\!\SI{0.40}{tCO_2/\mega\watt\hour}$),
\eqref{eq:switching-price} yields
$p^\star\!\approx\!\SI{56}{\dollar/tCO_2}$, a value squarely inside
the EU~ETS~2025 trading range.

\paragraph{(d) Practical reading.}
In the \cenpipe{} reconstruction, the dual pair
\eqref{eq:cost-primal}--\eqref{eq:cost-dual} is \emph{not} solved in
either direction.  The published $\lambda_t^{\mathrm{CEN}}$ is taken as
the realised optimal dual on \eqref{eq:cost-balance} and the
closest-CV picker recovers $m(t)$ from the data.  Once $m(t)$ is
identified, the emission-marginal value follows directly:
\begin{equation}
  \hat\varepsilon_t \;=\; e_{\,\mathrm{fuel}(m(t))},
  \label{eq:eps-from-m}
\end{equation}
under the sufficient assumption that the emission-marginal unit
coincides with the cost-marginal unit, i.e.\ $m(t)\!=\!m'(t)$.  This
assumption holds whenever the marginal unit is the only generator
strictly interior to its bounds at block $t$ and is what justifies
treating $\hat\varepsilon_t$ as the true emission-marginal rate
without re-solving an emission-LP.  Two regimes break the
assumption: (i)~degenerate cost ties between two fuels (e.g.\ two gas
units of identical CV but different heat-rate-derived $e$), in which
case~\eqref{eq:eps-from-m} returns the $e_g$ of an arbitrary
representative; and (ii)~must-run units in their interior, which by
construction the picker excludes from the candidate pool but which
would otherwise compete for the emission margin.  Both pathologies
are flagged in the
\texttt{tier\_uncertainty} column of the output parquet and trigger a
\texttt{discrepancy\_flag} in the audit oracle.

\paragraph{Per-bus extension.}
The transmission-aware per-bus version of
\eqref{eq:cost-primal}--\eqref{eq:cost-dual} replaces the
single-bus balance with a vector balance plus DC power-flow
constraints.  The per-bus dual variable is
\begin{equation}
  \lambda_b
  \;=\; \lambda_{\mathrm{ref}} \;+\; \mathrm{LMC}_b \;+\; \mathrm{LCC}_b,
  \label{eq:lmp-decomp-cited}
\end{equation}
in the standard Hogan~\cite{Hogan1992} decomposition, where
$\mathrm{LMC}_b$ is the marginal-loss component and
$\mathrm{LCC}_b$ the marginal-congestion component.  The
\cenpipe{} pipeline absorbs $\mathrm{LMC}_b$ into the
empirical or PID-derived loss factor $\mathrm{MLF}_b$ and
$\mathrm{LCC}_b$ into the island partition itself: two buses are
placed in the same island if and only if their $\lambda$ differs
by less than the tolerance and the cumulative spread along the
chain remains under the cap, conditions that are tantamount to
$\mathrm{LCC}_b\!=\!\mathrm{LCC}_{b'}$ to within the tolerance.
The per-bus emission factor then inherits the same decomposition
\begin{equation}
  \varepsilon_b \;=\; \varepsilon_{\mathrm{ref}}
                    \;+\; \mathrm{EMC}_b
                    \;+\; \mathrm{ECC}_b,
  \label{eq:eps-decomp}
\end{equation}
where $\mathrm{EMC}_b$ is the emission-marginal-loss component
($\varepsilon_{\mathrm{ref}}\!\cdot\!(\mathrm{MLF}_b\!-\!1)$ to first
order) and $\mathrm{ECC}_b$ vanishes inside any coherent island by
construction.  Equation~\eqref{eq:eps-decomp} is the parallel of
Hogan's LMP decomposition in the emission domain, and is, to the
best of our knowledge, the only published formulation that does so
without recourse to implicit-differentiation
co-simulation~\cite{ValenzuelaASL2022}.
